% Official source: https://www.imo-official.org/assets/documents/problems/2014/2014_eng.pdf
\competitionyear{2014}

\begin{problem}{2014}{1}
Let $a_0<a_1<a_2<\cdots$ be an infinite sequence of positive integers. Prove that there is exactly one integer $n\ge1$ for which
\[
    a_n<\frac{a_0+a_1+\cdots+a_n}{n}\le a_{n+1}.
\]
\end{problem}

\begin{problem}{2014}{2}
Let $n\ge2$ be an integer. A placement of $n$ rooks on an $n\times n$ chessboard is called \emph{peaceful} if each row and each column contains exactly one rook. Determine the largest positive integer $k$ such that every peaceful placement contains a $k\times k$ square of unit cells with no rook in any of its cells.
\end{problem}

\begin{problem}{2014}{3}
Let $ABCD$ be a convex quadrilateral with $\angle ABC=\angle CDA=90^\circ$. Let $H$ be the foot of the perpendicular from $A$ to $BD$. Points $S$ and $T$ lie on the sides $AB$ and $AD$, respectively, with $H$ inside triangle $SCT$ and
\[
    \angle CHS-\angle CSB=90^\circ,
    \qquad \angle THC-\angle DTC=90^\circ.
\]
Prove that the line $BD$ is tangent to the circumcircle of triangle $TSH$.
\end{problem}

\begin{problem}{2014}{4}
Let $P$ and $Q$ lie on the side $BC$ of an acute triangle $ABC$, with $\angle PAB=\angle BCA$ and $\angle CAQ=\angle ABC$. Choose $M$ on the line $AP$ and $N$ on the line $AQ$ so that $P$ and $Q$ are the midpoints of $AM$ and $AN$, respectively. Prove that the lines $BM$ and $CN$ intersect on the circumcircle of $ABC$.
\end{problem}

\begin{problem}{2014}{5}
For every positive integer $n$, the Bank of Cape Town issues coins worth $1/n$. Consider any finite collection of these coins, with denominations allowed to repeat, whose total value is at most $99+\frac12$. Prove that the coins can be divided into at most $100$ groups, each having total value at most $1$.
\end{problem}

\begin{problem}{2014}{6}
A set of lines in the plane is in \emph{general position} if no two are parallel and no three pass through one point. The lines divide the plane into regions; those with finite area are called \emph{finite regions}. Prove that, for all sufficiently large $n$, any $n$ lines in general position have at least $\sqrt n$ lines that can be coloured blue so that no finite region has a boundary made entirely of blue lines.
\end{problem}
